\documentclass[11pt,a4paper]{article}

% Packages
\usepackage[utf8]{inputenc}
\usepackage[T1]{fontenc}
\usepackage[margin=0.75in]{geometry}
\usepackage{enumitem}
\usepackage{hyperref}
\usepackage{titlesec}
\usepackage{xcolor}

% Formatting
\pagestyle{empty}
\setlength{\parindent}{0pt}
\setlist[itemize]{leftmargin=*, noitemsep, topsep=3pt}

% Section formatting
\titleformat{\section}{\Large\bfseries}{}{0em}{}[\titlerule]
\titlespacing*{\section}{0pt}{12pt}{6pt}

\titleformat{\subsection}[runin]{\bfseries}{}{0em}{}
\titlespacing*{\subsection}{0pt}{8pt}{0pt}

% Colors
\definecolor{linkcolor}{RGB}{0,102,204}
\hypersetup{
    colorlinks=true,
    urlcolor=linkcolor
}

\begin{document}

% Header
\begin{center}
    {\LARGE \textbf{ERIC MAYEUX}}\\
    \vspace{4pt}
    {\large Chargé de projet, Transformation Opérationnelle | Operational Transformation PM}\\
    \vspace{4pt}
    Montréal, QC, Canada\\
    \href{mailto:mayeux.e@gmail.com}{mayeux.e@gmail.com} | 438-884-7645
\end{center}

\vspace{8pt}

\section*{Profil Professionnel}
\noindent
\textbf{Chargé de projet en transformation opérationnelle} avec 10+ ans d'expérience en \textbf{déploiement de processus} et pilotage d'initiatives de changement en environnement TI. Expert en \textbf{gestion du changement}, planification d'ateliers et \textbf{cartographie de processus}. Maîtrise de la \textbf{suite Microsoft Office 365} et des tableaux de bord (Power BI). Certifié \textbf{PMP} (en cours), CSM. Autonome, rigoureux, bilinguisme français/anglais.

\section*{Compétences Techniques}

\textbf{Transformation \& Gestion du changement:} Déploiement de processus, plans de communication, gestion du changement, mitigation des risques, coaching des parties prenantes, amélioration continue

\textbf{Planification \& Suivi:} Indicateurs de mesure, tableaux de bord (Power BI), systèmes de contrôle, rapports d'avancement, prise de décision

\textbf{Animation \& Formation:} Planification d'ateliers, cartographie de processus, matériel de formation, facilitation de séances, animation de rencontres hebdomadaires

\textbf{Outils:} Microsoft Office 365, Power BI, Jira, Confluence, Datadog, Splunk

\textbf{Méthodologies:} Agile, Scrum, SAFe, Kanban, Waterfall, PMP

\textbf{Intelligence Artificielle:} Agents IA, automatisation de processus, RPA

\textbf{Architecture \& Cloud:} AWS Cloud Practitioner, Architectures Distribuées

\vspace{6pt}

\section*{Réalisations}
\begin{itemize}
    \item \textbf{Déploiement de processus} standards dans un contexte multi-équipes et multi-sites (BNC, Thaïlande)
    \item Élaboration de \textbf{plans de communication et gestion du changement} pour des transformations organisationnelles
    \item Mise en place d'\textbf{indicateurs de mesure} et tableaux de bord pour la prise de décision et le suivi de performance
    \item Animation d'\textbf{ateliers de travail} et de formations pour ancrer les meilleures pratiques
    \item \textbf{Coaching} d'équipes multidisciplinaires en contexte de transformation agile SAFe
    \item Organisation de hackathons autour de l'IA, de la migration et de l'\textbf{amélioration continue}
\end{itemize}

\section*{Expérience Professionnelle}

\subsection*{Scrum Master -- Gestion de projet | Project Manager} \hfill \textit{Novembre 2025 -- Présent}\\
\textbf{Banque Nationale du Canada (BNC)} -- Montréal, QC\\
Pilote la \textbf{transformation agile} et le déploiement de processus qualité dans un contexte bancaire.
\begin{itemize}
    \item Élabore et déploie des \textbf{plans de communication} et de gestion du changement pour les équipes TI
    \item Définit des \textbf{indicateurs de mesure} et systèmes de contrôle pour le suivi des initiatives
    \item \textbf{Planifie et anime des ateliers} de travail et des formations pour les équipes
    \item \textbf{Coache} les parties prenantes pour ancrer les pratiques agiles et l'amélioration continue
    \item Coordonne les enjeux inter-équipes ; produit des rapports d'avancement pour la haute direction
    \item Exploite \textbf{Microsoft Office 365} et des agents IA pour l'automatisation des tâches administratives
\end{itemize}

\subsection*{Intégrateur Sénior Qualité | Senior Quality Integrator} \hfill \textit{Octobre 2023 -- Novembre 2025}\\
\textbf{Banque Nationale du Canada (BNC)} -- Montréal, QC
\begin{itemize}
    \item \textbf{Cartographie des processus} qualité et déploiement de stratégies pan-train multi-sites
    \item Création de \textbf{tableaux de bord} (Datadog, Splunk, Jira) et mise en place de métriques qualité
    \item \textbf{Coaching} et formation continue des équipes QA ; organisation d'un hackathon de test
\end{itemize}

\subsection*{Intégrateur Qualité -- SDET | Quality Integrator -- Software Development Engineer in Test} \hfill \textit{Janvier 2022 -- Octobre 2023}\\
\textbf{Banque Nationale du Canada (BNC)} via Groupe Astek -- Montréal, QC
\begin{itemize}
    \item Mise en place de métriques qualité et \textbf{amélioration continue} des processus de test
    \item Stack : Selenium, AppliTools, RestAssured, JMeter, Gatling, K6
\end{itemize}

\subsection*{QA Automaticien | QA Automation Engineer} \hfill \textit{Juin 2021 -- Décembre 2021}\\
\textbf{AODocs} -- Paris, France
\begin{itemize}
    \item Automatisation des tests ; POC Flutter ; rapports Allure
\end{itemize}

\subsection*{Automaticien CI/CD | DevOps Automation Engineer} \hfill \textit{Janvier 2020 -- Juin 2021}\\
\textbf{ENGIE DIGITAL} via Groupe Hénix -- Paris, France
\begin{itemize}
    \item Création de la chaîne CI/CD : Jenkins, Docker ; pilotage projets Jira, KPIs
    \item Scripting Python, Shell, Bash pour l'\textbf{automatisation de processus}
\end{itemize}

\subsection*{Product Owner} \hfill \textit{Juin 2019 -- Septembre 2019}\\
\textbf{Hénix} -- Montrouge, France
\begin{itemize}
    \item Gestion du changement produit, backlog, support et formation clients SaaS
\end{itemize}

\subsection*{Développeur Logiciel | Software Developer} \hfill \textit{Mai 2017 -- Mai 2019}\\
\textbf{MEDIAPOST} via Groupe Hénix -- Paris, France
\begin{itemize}
    \item Développement back office ; reporting et macros Excel pour la direction
\end{itemize}

\subsection*{Consultant QA \& Automatisation | QA Automation Consultant} \hfill \textit{Avril 2016 -- Avril 2017}\\
\textbf{ENGIE IT} via Groupe Hénix -- Saint-Ouen, France
\begin{itemize}
    \item Pilotage opérationnel, gestion des incidents, administration de plateforme
\end{itemize}

\section*{Certifications \& Formations}

\textbf{PMP Certification} (en cours) \hfill \textit{2026}

Project Management Professional -- Project Management Institute (PMI)

\textbf{AWS Cloud Practitioner} \hfill \textit{2026}

Amazon Web Services (AWS)

\textbf{Certified Scrum Master (CSM)} \hfill \textit{2024}

Scrum Alliance

\textbf{Jira ACP-600 -- Project Administration in Jira Server} \hfill \textit{2019}

Atlassian Certified Professional

\textbf{Cursus Automaticien et DevOps} \hfill \textit{2019}\\
\textit{AFCEPF -- Montrouge, France}

\textbf{Cursus Architecture logiciel \& Analyste informaticien - Certifié Master II} \hfill \textit{2017}\\
\textit{AFCEPF-HENIX -- Malakoff, France}

\section*{Formation Académique}

\textbf{Licence en Gestion (Bachelor's Degree in Management)} \hfill \textit{2011}\\
École Supérieure de Gestion (E.S.G) -- Paris, France

\section*{Compétences Linguistiques}

\textbf{Français:} Langue maternelle (Native)\\
\textbf{Anglais:} Niveau Avancé (Advanced / Professional Working Proficiency)

\end{document}
